\documentclass[11pt,a4paper]{article}

% ─── Packages ───
\usepackage[utf8]{inputenc}
\usepackage[T1]{fontenc}
\usepackage{amsmath,amssymb,amsthm}
\usepackage{mathtools}
\usepackage{bm}
\usepackage{tikz}
\usepackage{booktabs}
\usepackage{hyperref}
\usepackage{cleveref}
\usepackage[margin=1in]{geometry}
\usepackage{enumitem}
\usepackage[expansion=false]{microtype}
\usepackage{xcolor}

% ─── Theorem environments ───
\newtheorem{definition}{Definition}
\newtheorem{proposition}{Proposition}
\newtheorem{theorem}{Theorem}
\newtheorem{remark}{Remark}
\newtheorem{corollary}{Corollary}

% ─── Notation ───
\newcommand{\B}{\mathcal{B}}
\newcommand{\X}{\mathcal{X}}
\newcommand{\R}{\mathbb{R}}
\newcommand{\Z}{\mathbb{Z}}
\newcommand{\N}{\mathbb{N}}
\newcommand{\bvec}{\bm{b}}
\newcommand{\xvec}{\bm{x}}
\newcommand{\yvec}{\bm{y}}
\newcommand{\dvec}{\bm{d}}
\newcommand{\phivec}{\bm{\phi}}

\title{%
  \textbf{Generalized Braille}:\\
  From Tactile Code to Inspectable\\
  Semantic Vector Space%
}
\author{
  Ryan Barrett \\
  Elevate Foundry \\
  \texttt{ryan@elevatefoundry.com}
}
\date{\today}

\begin{document}
\maketitle

\begin{abstract}
We formalize Braille as a family of structured codes over signed,
extensible, finite-dimensional vector spaces.
Standard 6-dot Braille is the instance $\B_6 = \{0,1\}^6$;
8-dot computer Braille is $\B_8 = \{0,1\}^8$.
We generalize to $n$-dot Braille over the ternary alphabet $\{-1,0,+1\}^n$,
which introduces explicit negation, inhibition, and algebraic structure
absent from binary encodings.
We show that composing ternary cells into sequences yields a
representational system with no fixed semantic-complexity ceiling,
and that the limiting factor shifts from vocabulary size to the
quality of the learned composition law.
We further extend to unbounded-dot and infinity-dot Braille,
connecting the framework to sparse feature spaces and
neural latent representations.
We describe a concrete implementation as a pre-training objective
for convolutional style-transfer networks, and present curriculum
experiments spanning Nemeth mathematical Braille, signed-dot patterns,
and $n$-dot cells for $n \in \{6, 8, 10, 12\}$.
\end{abstract}

% =====================================================================
\section{Introduction}
\label{sec:intro}
% =====================================================================

Braille is conventionally understood as a tactile writing system:
a fixed grid of raised dots encoding characters from a finite alphabet.
Standard Braille uses a $2 \times 3$ grid of 6 dots, yielding $2^6 = 64$
possible cell patterns.
Computer Braille extends this to a $2 \times 4$ grid of 8 dots,
yielding $2^8 = 256$ patterns.

In both cases, a cell is a binary vector:
\begin{equation}
  \bvec = (b_1, \ldots, b_n), \qquad b_i \in \{0,1\},
\end{equation}
and the set of all cells is:
\begin{equation}
  |\B_n| = 2^n.
\end{equation}

We observe that this is an instance of a much richer family of
structured codes.
The natural generalization replaces the binary alphabet with a
signed ternary alphabet, extends the dimensionality to arbitrary $n$,
and considers composition of cells into sequences.
The result is a framework that bridges tactile encoding,
symbolic representation, and neural feature spaces.

% =====================================================================
\section{The \texorpdfstring{$n$}{n}-dot Braille Family}
\label{sec:ndot}
% =====================================================================

\begin{definition}[$n$-dot binary Braille]
  An $n$-dot binary Braille cell is a vector
  $\bvec \in \{0,1\}^n$,
  arranged in a $2 \times \lceil n/2 \rceil$ spatial grid.
  The cell space has cardinality $2^n$.
\end{definition}

The standard instances form a nested hierarchy:
\begin{equation}
  \B_6 \subset \B_8 \subset \B_{10} \subset \cdots
\end{equation}
assuming the added dimensions preserve the semantics of earlier cells
(i.e., dots 1--6 retain their standard meanings).

\begin{table}[h]
  \centering
  \begin{tabular}{@{}lccc@{}}
    \toprule
    System & $n$ & Grid & $|\B_n|$ \\
    \midrule
    Standard Braille & 6 & $2 \times 3$ & 64 \\
    Computer Braille & 8 & $2 \times 4$ & 256 \\
    Extended Braille & 10 & $2 \times 5$ & 1{,}024 \\
    12-dot Braille & 12 & $2 \times 6$ & 4{,}096 \\
    \bottomrule
  \end{tabular}
  \caption{Binary $n$-dot Braille instances.}
  \label{tab:ndot}
\end{table}

The spatial arrangement matters.
A cell is not merely a bit string; it is a \emph{spatially grounded}
bit string.
The 2-column structure assigns each dot a position $(c, r)$ where
$c \in \{0, 1\}$ (left/right) and $r \in \{0, \ldots, \lceil n/2 \rceil - 1\}$ (row).
This spatial grounding is what distinguishes Braille from an
arbitrary binary code and makes it learnable by spatial perception systems.

% =====================================================================
\section{Signed Braille: The Ternary Extension}
\label{sec:signed}
% =====================================================================

\begin{definition}[Signed $n$-dot Braille]
  A signed Braille cell is a vector
  $\bvec \in \{-1, 0, +1\}^n$,
  where:
  \begin{align}
    b_i = +1 &\quad \text{(asserted: the feature is present)}, \\
    b_i = 0  &\quad \text{(unspecified: no statement about feature $i$)}, \\
    b_i = -1 &\quad \text{(denied: the feature is explicitly absent)}.
  \end{align}
\end{definition}

The cell space has cardinality $3^n$.
For 8-dot signed Braille: $3^8 = 6{,}561$ patterns.
For 10-dot: $3^{10} = 59{,}049$.

\begin{remark}
  Binary Braille cannot distinguish ``not present'' from
  ``explicitly absent.''
  The value $b_i = 0$ in $\{0,1\}^n$ conflates two distinct semantic states.
  Signed Braille resolves this conflation.
\end{remark}

\subsection{Algebraic Properties}

Signed cells admit natural algebraic operations.

\paragraph{Negation.}
For $\xvec = (x_1, \ldots, x_n)$:
\begin{equation}
  -\xvec = (-x_1, \ldots, -x_n).
\end{equation}
This corresponds to conceptual inversion: every asserted feature
becomes denied, and vice versa.

\paragraph{Cancellation.}
\begin{equation}
  (+d_i) + (-d_i) = 0.
\end{equation}
A denied dot cancels an asserted dot, yielding the unspecified state.

\paragraph{Inner product.}
\begin{equation}
  \langle \xvec, \yvec \rangle = \sum_{i=1}^{n} x_i y_i.
\end{equation}
This measures alignment versus opposition:
\begin{equation}
  x_i y_i =
  \begin{cases}
    +1 & \text{agreement (both asserted or both denied)}, \\
    0  & \text{irrelevant (at least one unspecified)}, \\
    -1 & \text{contradiction (one asserted, one denied)}.
  \end{cases}
\end{equation}

\paragraph{Signed similarity.}
Two cells agree when their signs align and conflict when they oppose.
This permits distance metrics, clustering, and analogy operations
directly on the cell representation.

% =====================================================================
\section{Compositional Expressivity}
\label{sec:composition}
% =====================================================================

A single $n$-dot ternary cell carries:
\begin{equation}
  \log_2(3^n) = n \log_2 3 \approx 1.585n \text{ bits}.
\end{equation}

For 10-dot ternary Braille, this is approximately 15.85 bits per cell.
Substantial, but finite.

The real expressivity arises from \emph{composition}.
A sequence of $L$ cells:
\begin{equation}
  \xvec = (\xvec_1, \xvec_2, \ldots, \xvec_L),
\end{equation}
lives in the product space:
\begin{equation}
  |\X_L| = (3^n)^L = 3^{nL}.
\end{equation}

At only 32 cells of 10-dot ternary Braille:
\begin{equation}
  3^{320} \approx 10^{153}.
\end{equation}

But state counting misses the point.
English has fewer than 100 primitive characters, yet:
\begin{equation}
  \text{small alphabet} + \text{composition}
  \rightarrow
  \text{Shakespeare, mathematics, Python}.
\end{equation}

Binary computers are even more extreme: $\{0,1\}$ suffices for
essentially all digitally representable knowledge.

\begin{theorem}[Expressivity hierarchy]
  Ternary $n$-dot Braille with composition admits three levels
  of expressivity:
  \begin{enumerate}[label=\arabic*.]
    \item \textbf{Atomic:} $3^n$ meanings (finite, fixed).
    \item \textbf{Linguistic:} $(\{-1,0,+1\}^n)^*$ ---
      infinitely many finite expressions encoding
      sentences, trees, graphs, proofs, and programs.
    \item \textbf{Computational:} With read, write, branch, bind,
      and recursion operations plus unbounded working memory,
      the system is Turing-complete.
  \end{enumerate}
\end{theorem}

Therefore:
\begin{equation}
  \boxed{
    \{-1,0,+1\}^n + \text{composition}
    \rightarrow
    \text{arbitrary finite symbolic structures}
  }
\end{equation}
and:
\begin{equation}
  \boxed{
    \{-1,0,+1\}^n + \text{composition} + \text{recursion} + \text{memory}
    \rightarrow
    \text{general computation}
  }
\end{equation}

% =====================================================================
\section{Factored Representations}
\label{sec:factored}
% =====================================================================

The crucial design principle: meaning should \emph{factor through the dots}.

An arbitrary lookup table assigns:
\begin{equation}
  \bvec \mapsto \text{symbol}.
\end{equation}

A semantic Braille system assigns:
\begin{equation}
  b_i \mapsto \text{primitive semantic feature}
\end{equation}
and therefore:
\begin{equation}
  \bvec \mapsto \text{composition of active features}.
\end{equation}

Increasing $n$ then adds \emph{semantic axes} rather than merely
vocabulary capacity.

The cell-level embedding becomes:
\begin{equation}
  \phi(\bvec) = \sum_{i=1}^{n} b_i \dvec_i
    + \sum_{i < j} b_i b_j \dvec_{ij} + \cdots
  \label{eq:factored}
\end{equation}
where $\dvec_i$ represents a primitive dot meaning and
$\dvec_{ij}$ represents pairwise interactions.

Sequence-level composition adds another layer:
\begin{equation}
  \Phi(\xvec_1, \ldots, \xvec_L) = F(\phi(\xvec_1), \ldots, \phi(\xvec_L)).
\end{equation}

The full hierarchy is:
\begin{equation}
  \text{trit}
  \rightarrow \text{dot interactions}
  \rightarrow \text{cell}
  \rightarrow \text{cell interactions}
  \rightarrow \text{structure}
  \rightarrow \text{program/meaning}.
\end{equation}

\begin{remark}
  Past some $n$, exhaustive learning of all $3^n$ cells becomes
  impossible.
  The system must learn the \emph{factorization of dots}
  (Equation~\ref{eq:factored}), not every possible cell independently.
  This is the regime where factored representations become essential.
\end{remark}

% =====================================================================
\section{Extensions Beyond Ternary}
\label{sec:extensions}
% =====================================================================

The ternary alphabet is not the endpoint.

\begin{definition}[Generalized dot algebra]
  A dot may take values in an arbitrary domain $D_i$:
  \begin{align}
    b_i \in \Z       &\quad \text{(integer: count, multiplicity)}, \\
    b_i \in \R       &\quad \text{(real: activation, confidence)}, \\
    b_i \in \mathbb{C} &\quad \text{(complex: magnitude + phase)}, \\
    b_i \in [0,1]    &\quad \text{(probabilistic: belief)}, \\
    b_i \in [l_i, u_i] &\quad \text{(interval: uncertainty bounds)}.
  \end{align}
\end{definition}

The most coherent hierarchy is:
\begin{align}
  \text{ordinary Braille}   &= \{0,1\}^n, \\
  \text{signed Braille}     &= \{-1,0,+1\}^n, \\
  \text{weighted Braille}   &= \R^n, \\
  \text{infinity-dot Braille} &= \R^{\N} \text{ with sparse or computable support}.
\end{align}

This converges toward a learned latent space, but with a critical
difference: \textbf{the dimensions remain externally addressable
and potentially inspectable}.

A neural vector is usually $\bm{z} \in \R^d$ with opaque coordinates.
Generalized Braille is $\bvec \in \R^d$ with spatially or semantically
grounded coordinates.

\begin{equation}
  \boxed{
    \text{Generalized Braille}
    = \text{an inspectable, extensible, signed semantic vector space}
  }
\end{equation}

% =====================================================================
\section{Infinity-dot Braille}
\label{sec:infinity}
% =====================================================================

\subsection{Unbounded Braille}

The cell may contain any finite number of dimensions, with no fixed
maximum:
\begin{equation}
  \bvec \in \bigcup_{n=1}^{\infty} \{0,1\}^n.
\end{equation}

This is operationally sensible.
A representation can expand when it needs new semantic dimensions.
It resembles a dynamically extensible schema:
\begin{equation}
  \bvec = (b_1, \ldots, b_n), \qquad n \text{ chosen per expression}.
\end{equation}

This is the useful form of ``infinity-dot Braille.''

\subsection{Actually Infinite Braille}

A cell is an infinite binary sequence:
\begin{equation}
  \bvec \in \{0,1\}^{\N}.
\end{equation}

The complete space is uncountable:
\begin{equation}
  \left|\{0,1\}^{\N}\right| = 2^{\aleph_0}.
\end{equation}

But computation can only access finite prefixes.
Practical infinity-dot Braille must impose one of:
\begin{itemize}
  \item Finite support: $\sum_i b_i < \infty$.
  \item Computable sequence: $b_i = f(i)$ for computable $f$.
  \item Lazy representation: $b_{1:k}$ generated on demand.
\end{itemize}

A sparse infinite cell is especially natural:
\begin{equation}
  \bvec: \N \rightarrow \{0,1\}, \qquad
  |\{i : b_i = 1\}| < \infty.
\end{equation}

This gives an unlimited namespace while keeping every actual
expression finite:
\begin{equation}
  \bvec \equiv \{i_1, i_2, \ldots, i_k\}.
\end{equation}

\begin{equation}
  \boxed{
    \text{Infinity-dot Braille is a semantic coordinate system.}
  }
\end{equation}

% =====================================================================
\section{Implementation: Pre-training Objective}
\label{sec:implementation}
% =====================================================================

We implement generalized Braille as a pre-training objective for
convolutional style-transfer networks.

\subsection{Synthetic Data Generation}

For $n$-dot signed Braille with $n$ configurable:
\begin{enumerate}
  \item Generate all $3^n$ (or $2^n$) possible cell patterns.
  \item Render cells as images on a $2 \times \lceil n/2 \rceil$ spatial grid.
  \item Asserted dots ($+1$) are rendered as filled circles.
  \item Denied dots ($-1$) are rendered as X marks.
  \item Absent dots ($0$) produce no mark.
  \item Add position jitter, size variation, and pressure simulation
    to approximate hand-drawn appearance.
\end{enumerate}

\subsection{Network Architecture}

The pre-training target is a convolutional encoder-decoder
(TransformerNet) trained with perceptual loss via VGG-16 features.
The hypothesis: teaching the network that spatial dot patterns carry
structured, compositionally meaningful information produces better
spatial preservation during style transfer.

\subsection{Training Protocol}

\begin{enumerate}
  \item \textbf{Stage 1 (Braille):}
    Pre-train on synthetic signed $n$-dot Braille images.
    Content weight is doubled to emphasize spatial fidelity.
  \item \textbf{Stage 2 (Style transfer):}
    Fine-tune on natural images (COCO) with standard style loss.
\end{enumerate}

\subsection{Curriculum}

The drawing curriculum includes:
\begin{itemize}
  \item 8-dot lowercase and uppercase alphabet
    (dot 7 as capital indicator).
  \item Nemeth mathematical Braille: integrals, summations,
    Greek letters, the quadratic formula, Euler's identity,
    the Pythagorean theorem, and the Basel problem.
  \item Generalized $n$-dot hierarchy:
    $\B_6 \subset \B_8 \subset \B_{10}$ rendered side-by-side.
  \item Negative-dot demonstrations:
    asserted, denied, and absent in the same cell.
  \item 12-dot cells ($2 \times 6$ grid, $4{,}096$ binary patterns).
  \item Algebraic cancellation: $(+d) + (-d) = 0$.
\end{itemize}

% =====================================================================
\section{Conclusion}
\label{sec:conclusion}
% =====================================================================

We have shown that Braille, properly generalized, is not a fixed
alphabet but a \emph{semantic vector space} with the following
properties:

\begin{enumerate}
  \item \textbf{Spatial grounding:}
    Dot positions are physically meaningful, not arbitrary indices.
  \item \textbf{Signed semantics:}
    The ternary alphabet $\{-1, 0, +1\}$ distinguishes
    assertion, absence, and denial.
  \item \textbf{Compositionality:}
    Cell sequences can encode arbitrary finite structures.
  \item \textbf{Extensibility:}
    Adding dimensions adds semantic axes, not just capacity.
  \item \textbf{Inspectability:}
    Unlike opaque neural embeddings, dot coordinates are
    externally addressable.
  \item \textbf{Algebraic structure:}
    Negation, cancellation, and inner products operate
    directly on the representation.
\end{enumerate}

The key results:

\medskip

\begin{center}
\begin{tabular}{@{}ll@{}}
  \toprule
  Finite $n$-dot Braille & is a vocabulary \\
  Unbounded-dot Braille & is a growing ontology \\
  Infinity-dot Braille & is a semantic basis \\
  Negative-dot Braille & adds contradiction and inhibition \\
  \bottomrule
\end{tabular}
\end{center}

\medskip

A fixed ternary $n$-dot cell is finite.
But an arbitrarily composable ternary Braille language has
\emph{no fixed semantic-complexity ceiling}.
The limiting factor is not the number of dots.
It is the quality of the learned composition law.

\begin{equation}
  \boxed{
    \text{Generalized Braille}
    = \text{an inspectable, extensible, signed semantic vector space.}
  }
\end{equation}

\bibliographystyle{plain}
% \bibliography{references}

\end{document}
